% Algebra 2 Lesson 4-1 fast-track packet: Inverse Variation and the Reciprocal Function (enVision A2, Example 3 only; brief lesson)
% Keep/cut is a proposal (2026-10-30); the PLC has not met on this lesson.
% Example 3 is kept because its two steps (find k from one known pair, then use it) answer quiz items 1-3,
% and it is the only Example with DOK 3 practice items (20, 22, 25, 26). Examples 4-5 (graphs, quiz 4-5) are cut.
% Shared preamble for the "fast" (hacked-down) Algebra 2 lesson packets, Oct 2026.
% Each packet: % Shared preamble for the "fast" (hacked-down) Algebra 2 lesson packets, Oct 2026.
% Each packet: % Shared preamble for the "fast" (hacked-down) Algebra 2 lesson packets, Oct 2026.
% Each packet: \input{fast-packet-template} then \begin{document} ... \end{document}.
\documentclass[12pt]{article}
\usepackage[margin=0.7in]{geometry}
\usepackage{amsmath,amssymb}
\usepackage{booktabs,tabularx,array}
\usepackage{enumitem}
\usepackage[T1]{fontenc}
\usepackage[utf8]{inputenc}
\usepackage{lmodern}
\usepackage{tikz}
\usepackage{pgfplots}
\pgfplotsset{compat=1.18}
\usepackage{xcolor}
\setlist{nosep,leftmargin=1.4em}
\setlength{\parindent}{0pt}
\setlength{\parskip}{6pt}
\renewcommand{\arraystretch}{1.35}
\newcolumntype{Y}{>{\raggedright\arraybackslash}X}
\pagestyle{empty}

% blanks and workspace
\newcommand{\blank}[1][2.6cm]{\rule[-2pt]{#1}{0.5pt}}
\newcommand{\lines}[1][1]{\par\vspace{2pt}\foreach \i in {1,...,#1}{\rule{0pt}{22pt}\rule{\linewidth}{0.4pt}\par}}
\newcommand{\wspart}[2]{\vspace{4pt}{\large\bfseries #1: #2}\par\vspace{-6pt}\rule{\textwidth}{0.4pt}\par}

% header: \packethead{Lesson 1-3}{Piecewise-Defined Functions}
\newcommand{\packethead}[2]{%
{\LARGE\bfseries #2}\hfill Name \blank[5cm]\\[2pt]
Algebra 2 \quad #1 (fast track) \hfill Date \blank[2.5cm] \quad Period \blank[1.2cm]\par\vspace{2pt}}

% boxed goals: \goals{learning goal}{language goal}
\newcommand{\goals}[2]{%
\fbox{\parbox{\dimexpr\textwidth-2\fboxsep-2\fboxrule}{%
\textbf{Learning goal:} #1\\[2pt]
\textbf{Language goal:} #2}}\par}

% quiz target box: \quiztarget{text}
\newcommand{\quiztarget}[1]{%
\vspace{4pt}\colorbox{gray!12}{\parbox{\dimexpr\textwidth-2\fboxsep}{\textbf{On the quiz:} #1}}\par}

% IXL table row helper: code & skill & which quiz question it feeds
\newcommand{\ixlrow}[3]{\texttt{#1} & #2 & #3 \\}

% one grid style for every graph; override xmin/xmax/... per plot
\pgfplotsset{sheetgrid/.style={
  axis lines=middle, grid=both, grid style={gray!45},
  xmin=-10.5, xmax=10.5, ymin=-10.5, ymax=10.5,
  xtick={-10,-8,...,10}, ytick={-10,-8,...,10},
  x=0.28cm, y=0.28cm,
  ticklabel style={font=\scriptsize, fill=white, inner sep=1pt},
  xlabel={$x$}, ylabel={$y$},
  axis line style={thick}, clip=true}}
% unit grid for small windows
\pgfplotsset{unitgrid/.style={sheetgrid,
  xmin=-5.5, xmax=5.5, ymin=-5.5, ymax=5.5,
  xtick={-5,...,5}, ytick={-5,...,5}, x=0.5cm, y=0.5cm}}
% open and closed endpoints
\tikzset{closed/.style={circle, fill, inner sep=1.6pt},
         open/.style={circle, draw, fill=white, inner sep=1.4pt}}
 then \begin{document} ... \end{document}.
\documentclass[12pt]{article}
\usepackage[margin=0.7in]{geometry}
\usepackage{amsmath,amssymb}
\usepackage{booktabs,tabularx,array}
\usepackage{enumitem}
\usepackage[T1]{fontenc}
\usepackage[utf8]{inputenc}
\usepackage{lmodern}
\usepackage{tikz}
\usepackage{pgfplots}
\pgfplotsset{compat=1.18}
\usepackage{xcolor}
\setlist{nosep,leftmargin=1.4em}
\setlength{\parindent}{0pt}
\setlength{\parskip}{6pt}
\renewcommand{\arraystretch}{1.35}
\newcolumntype{Y}{>{\raggedright\arraybackslash}X}
\pagestyle{empty}

% blanks and workspace
\newcommand{\blank}[1][2.6cm]{\rule[-2pt]{#1}{0.5pt}}
\newcommand{\lines}[1][1]{\par\vspace{2pt}\foreach \i in {1,...,#1}{\rule{0pt}{22pt}\rule{\linewidth}{0.4pt}\par}}
\newcommand{\wspart}[2]{\vspace{4pt}{\large\bfseries #1: #2}\par\vspace{-6pt}\rule{\textwidth}{0.4pt}\par}

% header: \packethead{Lesson 1-3}{Piecewise-Defined Functions}
\newcommand{\packethead}[2]{%
{\LARGE\bfseries #2}\hfill Name \blank[5cm]\\[2pt]
Algebra 2 \quad #1 (fast track) \hfill Date \blank[2.5cm] \quad Period \blank[1.2cm]\par\vspace{2pt}}

% boxed goals: \goals{learning goal}{language goal}
\newcommand{\goals}[2]{%
\fbox{\parbox{\dimexpr\textwidth-2\fboxsep-2\fboxrule}{%
\textbf{Learning goal:} #1\\[2pt]
\textbf{Language goal:} #2}}\par}

% quiz target box: \quiztarget{text}
\newcommand{\quiztarget}[1]{%
\vspace{4pt}\colorbox{gray!12}{\parbox{\dimexpr\textwidth-2\fboxsep}{\textbf{On the quiz:} #1}}\par}

% IXL table row helper: code & skill & which quiz question it feeds
\newcommand{\ixlrow}[3]{\texttt{#1} & #2 & #3 \\}

% one grid style for every graph; override xmin/xmax/... per plot
\pgfplotsset{sheetgrid/.style={
  axis lines=middle, grid=both, grid style={gray!45},
  xmin=-10.5, xmax=10.5, ymin=-10.5, ymax=10.5,
  xtick={-10,-8,...,10}, ytick={-10,-8,...,10},
  x=0.28cm, y=0.28cm,
  ticklabel style={font=\scriptsize, fill=white, inner sep=1pt},
  xlabel={$x$}, ylabel={$y$},
  axis line style={thick}, clip=true}}
% unit grid for small windows
\pgfplotsset{unitgrid/.style={sheetgrid,
  xmin=-5.5, xmax=5.5, ymin=-5.5, ymax=5.5,
  xtick={-5,...,5}, ytick={-5,...,5}, x=0.5cm, y=0.5cm}}
% open and closed endpoints
\tikzset{closed/.style={circle, fill, inner sep=1.6pt},
         open/.style={circle, draw, fill=white, inner sep=1.4pt}}
 then \begin{document} ... \end{document}.
\documentclass[12pt]{article}
\usepackage[margin=0.7in]{geometry}
\usepackage{amsmath,amssymb}
\usepackage{booktabs,tabularx,array}
\usepackage{enumitem}
\usepackage[T1]{fontenc}
\usepackage[utf8]{inputenc}
\usepackage{lmodern}
\usepackage{tikz}
\usepackage{pgfplots}
\pgfplotsset{compat=1.18}
\usepackage{xcolor}
\setlist{nosep,leftmargin=1.4em}
\setlength{\parindent}{0pt}
\setlength{\parskip}{6pt}
\renewcommand{\arraystretch}{1.35}
\newcolumntype{Y}{>{\raggedright\arraybackslash}X}
\pagestyle{empty}

% blanks and workspace
\newcommand{\blank}[1][2.6cm]{\rule[-2pt]{#1}{0.5pt}}
\newcommand{\lines}[1][1]{\par\vspace{2pt}\foreach \i in {1,...,#1}{\rule{0pt}{22pt}\rule{\linewidth}{0.4pt}\par}}
\newcommand{\wspart}[2]{\vspace{4pt}{\large\bfseries #1: #2}\par\vspace{-6pt}\rule{\textwidth}{0.4pt}\par}

% header: \packethead{Lesson 1-3}{Piecewise-Defined Functions}
\newcommand{\packethead}[2]{%
{\LARGE\bfseries #2}\hfill Name \blank[5cm]\\[2pt]
Algebra 2 \quad #1 (fast track) \hfill Date \blank[2.5cm] \quad Period \blank[1.2cm]\par\vspace{2pt}}

% boxed goals: \goals{learning goal}{language goal}
\newcommand{\goals}[2]{%
\fbox{\parbox{\dimexpr\textwidth-2\fboxsep-2\fboxrule}{%
\textbf{Learning goal:} #1\\[2pt]
\textbf{Language goal:} #2}}\par}

% quiz target box: \quiztarget{text}
\newcommand{\quiztarget}[1]{%
\vspace{4pt}\colorbox{gray!12}{\parbox{\dimexpr\textwidth-2\fboxsep}{\textbf{On the quiz:} #1}}\par}

% IXL table row helper: code & skill & which quiz question it feeds
\newcommand{\ixlrow}[3]{\texttt{#1} & #2 & #3 \\}

% one grid style for every graph; override xmin/xmax/... per plot
\pgfplotsset{sheetgrid/.style={
  axis lines=middle, grid=both, grid style={gray!45},
  xmin=-10.5, xmax=10.5, ymin=-10.5, ymax=10.5,
  xtick={-10,-8,...,10}, ytick={-10,-8,...,10},
  x=0.28cm, y=0.28cm,
  ticklabel style={font=\scriptsize, fill=white, inner sep=1pt},
  xlabel={$x$}, ylabel={$y$},
  axis line style={thick}, clip=true}}
% unit grid for small windows
\pgfplotsset{unitgrid/.style={sheetgrid,
  xmin=-5.5, xmax=5.5, ymin=-5.5, ymax=5.5,
  xtick={-5,...,5}, ytick={-5,...,5}, x=0.5cm, y=0.5cm}}
% open and closed endpoints
\tikzset{closed/.style={circle, fill, inner sep=1.6pt},
         open/.style={circle, draw, fill=white, inner sep=1.4pt}}

\begin{document}

\packethead{Lesson 4-1}{Inverse Variation and the Reciprocal Function}
\goals{I can write an inverse variation equation from one pair of values and use it to find a missing value in a table or a word problem.}{I can explain how I found a missing value using the words \emph{varies inversely}, \emph{product}, \emph{constant of variation}, and \emph{inverse variation}.}

% Spanish: the TE vocabulary builder prints "constant de proporcionalidad"; its \te note reads "likely constante" (corrected here).
{\footnotesize Spanish: \emph{variaci\'on inversa} (inverse variation), \emph{constante de proporcionalidad} (constant of variation), \emph{rec\'iproco} (reciprocal).}

\quiztarget{You will (1) find a missing value in an inverse-variation table, (2) choose the equation $y=\frac{k}{x}$ that fits a table, and (3) solve a word problem where one quantity varies inversely with another. Example 3 gives you the two steps for all three: find $k$ from a pair you know, then use $k$. Items (4) and (5), asymptotes and translations of $y=\frac{1}{x}$, are on the full Savvas quiz only.}

\vspace{2pt}
\fbox{\parbox{\dimexpr\textwidth-2\fboxsep-2\fboxrule}{%
\begin{minipage}[c]{0.45\linewidth}
\textbf{Key idea.} An \textbf{inverse variation} is modeled by $y=\frac{k}{x}$, where $k\neq0$, so $xy=k$. The \textbf{product} of every pair is the same number $k$, the \textbf{constant of variation}. For positive $x$ and $y$, as one increases, the other decreases.
% The TE's "as one variable increases, the other decreases proportionally" (concept box / concept summary WORDS; a2_4-1_TE.tex line 188 and 487) assumes k > 0 on one branch (TE \te note line 60); restated with its condition.
\end{minipage}\hfill
\begin{minipage}[c]{0.52\linewidth}
\small\setlength{\tabcolsep}{3.5pt}\renewcommand{\arraystretch}{1.15}
\begin{tabular}{|c|c|c|c|c|c|c|c|c|}
\hline
$x$ & 1 & 2 & 3 & 4 & 6 & 8 & 12 & 24 \\ \hline
$y$ & 24 & 12 & 8 & 6 & 4 & 3 & 2 & 1 \\ \hline
$k=xy$ & 24 & 24 & 24 & 24 & 24 & 24 & 24 & 24 \\ \hline
\end{tabular}

\vspace{2pt}
When $x$ doubles (1, 2, 4, 8), $y$ is halved (24, 12, 6, 3).
\end{minipage}}}

\wspart{Example 3}{Use an Inverse Variation Model}
On a Greek bouzouki, the string length $s$ varies inversely with the frequency $f$ of its vibrations. (Players can vary string length by placing their fingers behind the frets.) The frequency of a 26-inch E-string is 329.63 cycles per second. What is the frequency when the string length is 13 inches?

\textbf{Step 1} Write the equation for an inverse variation and solve for $k$ (26 for $s$, 329.63 for $f$).

\hspace*{1.5em}$s=\dfrac{k}{f}\qquad\longrightarrow\qquad 26=\dfrac{k}{329.63}\qquad\longrightarrow\qquad 26(329.63)=k\qquad\longrightarrow\qquad \mathbf{8{,}570.38=k}$

\textbf{Step 2} Substitute $k=8{,}570.38$ in the inverse variation equation and find $f$.

\hspace*{1.5em}$13=\dfrac{8{,}570.38}{f}\qquad\longrightarrow\qquad 13f=8{,}570.38\qquad\longrightarrow\qquad f=\dfrac{8{,}570.38}{13}=\mathbf{659.26}$

\begin{minipage}[c]{0.56\textwidth}
\textbf{Answer:} the 13-inch string has a frequency of \textbf{659.26 cycles per second}. \textbf{Check:} $13(659.26)=8{,}570.38$, the same product. \checkmark\ Half the length, double the frequency.
\end{minipage}\hfill
\begin{minipage}[c]{0.42\textwidth}
\small\renewcommand{\arraystretch}{1.2}
\hfill\begin{tabular}{|c|c|c|}
\hline
$s$ (in.) & $f$ (cycles/s) & $k=sf$ \\ \hline
26 & 329.63 & 8{,}570.38 \\ \hline
13 & \textbf{659.26} & 8{,}570.38 \\ \hline
\end{tabular}
\end{minipage}

\vspace{2pt}
\textbf{Say it:} The string length \emph{varies inversely} with the frequency, so the \emph{product} $sf$ is always \blank[2.2cm]. That number is the \emph{constant of variation}, $k$. When the length is cut in half, the frequency \blank[2cm]\\ because
\lines[2]

\newpage
\wspart{Try It 3}{Use an inverse variation model}
The amount of time it takes for an ice cube to melt varies inversely to the air temperature in degrees Celsius. At $20^\circ$ Celsius, the ice will melt in 20 minutes. How long will it take the ice to melt if the temperature is $30^\circ$ Celsius?

\begin{minipage}[t]{0.48\textwidth}
\textbf{Step 1} Let $t$ = minutes and $d$ = degrees Celsius. Write $t=\frac{k}{d}$, substitute the pair you know, and solve for $k$.

\vspace{4pt}
$k=t\cdot d=\blank[1.2cm]\cdot\blank[1.2cm]=\blank[1.8cm]$
\lines[2]
\end{minipage}\hfill
\begin{minipage}[t]{0.48\textwidth}
\textbf{Step 2} Substitute $k$ and the new temperature. Find $t$.

\vspace{4pt}
$t=\blank[1.6cm]\div\blank[1.2cm]\approx\blank[1.8cm]$ minutes
\lines[2]
\end{minipage}

\textbf{Check:} multiply your new pair, using the unrounded time (a fraction), not the rounded one. Is the product the same $k$? \blank[4cm]

Does it make sense that warmer air gives a \emph{shorter} time? Use the words \emph{varies inversely}.\lines[2]

% R5: quiz items 1 and 2 are tables (find a missing value; choose y = k/x with a negative k).
% Try It 3 is a word problem, so this Practice is a structural parallel of quiz items 1 and 2 (different numbers), not the items.
\wspart{Practice}{An inverse variation in a table}
The table represents an inverse variation.

\noindent\hfil{\renewcommand{\arraystretch}{1.5}\begin{tabular}{|c|*{4}{>{\centering\arraybackslash}p{1.5cm}|}}
\hline
$x$ & $-5$ & $-2$ & $4$ & $10$ \\ \hline
$y$ & $4$ & $10$ & $q$ & $-2$ \\ \hline
$xy$ & \blank[1.2cm] & \blank[1.2cm] & & \blank[1.2cm] \\ \hline
\end{tabular}}\hfil\par

\begin{enumerate}[label=\textbf{\alph*.}, itemsep=6pt]
\item Find the product $xy$ for each column you know. The constant of variation is $k=\blank[1.6cm]$.
\item Write the equation of the inverse variation: \quad $y=\blank[2.4cm]$
\item Find the value of $q$. Show your work.\lines[2]
\item Check: multiply $4$ and your $q$. Do you get $k$? \blank[3cm]

\vspace{2pt}
One of $x$, $y$ is negative in every pair because the \emph{product} $xy=k$ is \blank[3cm].
\end{enumerate}

\newpage
\wspart{IXL}{Practice for the quiz}
Go to IXL, Algebra 2. Type the code in the search box. Do the starred skills first.

\vspace{2pt}
\begin{tabularx}{\textwidth}{l Y Y}
\toprule
\textbf{Code} & \textbf{Skill} & \textbf{Feeds} \\
\midrule
\ixlrow{C9D}{Classify variation (warm-up)}{Warm-up: a constant product $xy$ means inverse variation}
\ixlrow{PNY $\bigstar$}{Write and solve inverse variation equations}{Quiz 1, 2 and 3: find $k$, then use it (Example 3)}
\ixlrow{PXE $\bigstar$}{Find the constant of variation}{Quiz 1 and 2: $k=xy$ from a table}
\ixlrow{ULD}{Domain and range of reciprocal functions (optional)}{Quiz 4 and 5 (full Savvas quiz only): asymptotes and translations of $y=1/(x-h)+k$; not in this packet}
\bottomrule
\end{tabularx}

\vspace{4pt}
\textbf{SmartScore 70 or above on each starred ($\bigstar$) skill replaces your quiz score.}

% R3: the lesson's DOK 3 item, Practice & Problem Solving 20 (Example 3, DOK 3), verbatim apart from "How long it will take" -> "How long will it take".
% TE key: t = 4,800/p; 24 min.
\wspart{Exit question}{Model with mathematics}
\fbox{\parbox{\dimexpr\textwidth-2\fboxsep-2\fboxrule}{%
The time $t$ required to empty a water tank varies inversely as the rate of pumping $p$. A pump can empty a water tank in 40 min at the rate of 120 gal/min. Write the equation of the inverse variation. How long will it take the pump to empty the water tank at the rate of 200 gal/min?

\vspace{4pt}
Equation: $t=\blank[3cm]$ \qquad Time at 200 gal/min: \blank[2.4cm] min
\lines[4]

\vspace{2pt}
Because $t$ \emph{varies inversely} as $p$, the \emph{product} \blank[1.6cm] is the \emph{constant of variation}, $k=\blank[2cm]$. A faster pump means \blank[4cm] time, because \dots
\lines[2]
}}

\vspace{6pt}
\textbf{Before the quiz, finish these sentences.}

In an \emph{inverse variation}, the \emph{product} $xy$ is \blank[4cm] for every pair.

To find a missing value, first find \blank[3cm], then \blank[5cm].

\end{document}
